In this chapter, we study some randomization techniques for faster
convex optimization.

\section{Subspace embedding\label{sec:Subspace-embedding}}

In this section and the next section, we consider the least squares
regression problem
\[
\min_{x}\|Ax-b\|^{2}_{2}
\]
where $A\in\R^{n\times d}$ with $n\ge d$ (we assume this throughout
the chapter). The gradient of the function is $2A^{\top}Ax-2A^{\top}b$.
Setting it to zero, and assuming $A^{\top}A$ is invertible, the solution
is given by 
\[
x=(A^{\top}A)^{-1}A^{\top}b.
\]
If $A^{\top}A$ is not invertible, we use its pseudo-inverse. If the
matrix $A^{\top}A\in\R^{d\times d}$ is given, then we can solve the
equation above in time $d^{\omega}$, the current complexity of matrix
multiplication. If $n>d^{\omega}$, then the bottleneck is simply
to compute $A^{\top}A$. The following lemma shows that it suffices
to approximate $A^{\top}A$. 

The simplest iteration is the Richardson iteration: 
\[
x^{(k+1)}=A^{\top}b+(I-A^{\top}A)x^{(k)}=x^{(k)}+A^{\top}b-A^{\top}Ax^{(k)}.
\]
To ensure this converges, we scale down $A^{\top}A$ by its largest
eigenvalue so that $A^{\top}A\prec I$. This gives a bound of $O(\kappa(A^{\top}A)\log(1/\epsilon))$
on the number of iterations to get $\epsilon$ error where $\kappa(A^{\top}A)=\lambda_{\max}(A^{\top}A)/\lambda_{\min}(A^{\top}A)$.
More generally, one can use \emph{pre-conditioning. }Recall that for
a vector $v$, the squared norm $\norm v^{2}_{M}$ is defined as $v^{\top}Mv$.
\begin{lem}
\label{lem:Richardson}Given a matrix $M$ such that $A^{\top}A\preceq M\preceq\kappa\cdot A^{\top}A$
for some $\kappa\geq1$. Consider the algorithm $x^{(k+1)}=x^{(k)}-M^{-1}(A^{\top}Ax^{(k)}-A^{\top}b)$.
Then, we have that
\[
\|x^{(k)}-x^{*}\|_{M}\leq\left(1-\frac{1}{\kappa}\right)^{k}\|x^{(0)}-x^{*}\|_{M}.
\]
\end{lem}

\begin{rem}
The proof also shows why the choice of norm above is the natural one.
In this norm, the residual drops geometrically.
\end{rem}

\begin{proof}
Using $x^{*}=(A^{\top}A)^{-1}A^{\top}b$, i.e., $A^{\top}b=(A^{\top}A)x^{*}$,
and the formula of $x^{(k+1)}$, we have
\begin{align*}
x^{(k+1)}-x^{*} & =x^{(k)}-M^{-1}(A^{\top}Ax^{(k)}-A^{\top}b)-x^{*}\\
 & =x^{(k)}-M^{-1}(A^{\top}Ax^{(k)}-A^{\top}Ax^{*})-x^{*}\\
 & =\left(I-M^{-1}A^{\top}A\right)\left(x^{(k)}-x^{*}\right).
\end{align*}
Therefore, the norm of error is
\begin{align*}
\|x^{(k+1)}-x^{*}\|^{2}_{M} & =\left(x^{(k)}-x^{*}\right)^{\top}(I-A^{\top}AM^{-1})M(I-M^{-1}A^{\top}A)\left(x^{(k)}-x^{*}\right)\\
 & =\left(x^{(k)}-x^{*}\right)^{\top}M^{\frac{1}{2}}(I-H)^{2}M^{\frac{1}{2}}\left(x^{(k)}-x^{*}\right)
\end{align*}
where $H=M^{-\frac{1}{2}}A^{\top}AM^{-\frac{1}{2}}$. Note that the
eigenvalues of $H$ lie between $1/\kappa$ and $1$ and hence 
\[
\lambda_{\max}(I-H)^{2}\leq\left(1-\frac{1}{\kappa}\right){}^{2}.
\]
This gives the conclusion.
\end{proof}
Note that $x^{\top}A^{\top}Ax=\|Ax\|^{2}$. Alternatively, we can
think of our goal as approximating $A$ by a smaller matrix $B$ such
that $\|Ax\|^{2}$ is close to $\|Bx\|^{2}$ for all $x$. In this
section, we show that we can simply take $B=\Pi A$ for a random matrix
$\Pi$ with relatively few rows. With this choice of $M$, we can
run the Richardson iteration. We need to see if this will make the
entire procedure more efficient.
\begin{defn}
A matrix $\Pi\in\R^{m\times n}$ is an embedding for a set $S\subset\R^{n}$
with distortion $\epsilon$ if 
\[
(1-\epsilon)\|y\|^{2}\leq\|\Pi y\|^{2}\leq(1+\epsilon)\|y\|^{2}
\]
for all $y\in S$.
\end{defn}

In this section, we focus on the case that $S$ is a $d$-dimensional
subspace in $\Rn$, namely $S=\{Ax:x\in\R^{d}\}$. Consider the SVD
$A=U\Sigma V^{\top}$. For any $y\in S$, we have that
\[
\|U^{\top}y\|^{2}_{2}=\|U^{\top}U\Sigma V^{\top}x\|^{2}_{2}=x^{\top}V\Sigma^{2}V^{\top}x=\|Ax\|^{2}.
\]
Therefore, any $d$-dimensional subspace has a zero distortion embedding
of size $d\times n$.
\begin{xca}
For any $d$-dimensional subspace $S$, any embedding with distortion
$\epsilon<1$ must have at least $d$ rows.
\end{xca}

This embedding is not useful for solving the least squares problem
because the solution of the least squares problem is simply a closed
form of the SVD decomposition $x=V\Sigma^{-1}U^{\top}b$ and finding
the SVD is usually more expensive.

\subsection{Oblivious Subspace Embedding via Johnson-Lindenstrauss}

Surprisingly, there are random embeddings that have small distortion
without knowledge of the subspace.
\begin{defn}
A random matrix $\Pi\in\R^{m\times n}$ is a $(d,\epsilon,\delta)$-oblivious
subspace embedding (OSE) if for any fixed $d$-dimensional subspace
$S\subset\R^{n}$, $\Pi$ is an embedding for $S$ with distortion
$\epsilon$ with probability at least $1-\delta$.
\end{defn}

The next lemma provides an equivalent definition.
\begin{lem}
We call $\Pi$ a $(d,\epsilon,\delta)$ OSE if for any matrix $U\in\R^{n\times d}$
with orthonormal columns, we have that
\[
\P(\|U^{\top}\Pi^{\top}\Pi U-I_{d}\|_{\op}\leq\epsilon)\geq1-\delta.
\]
\end{lem}

\begin{proof}
Let $S$ be the subspace with an orthonormal basis $U\in\R^{n\times d}$,
namely $S=\{y:y=Uz\}$. Then, the condition
\[
(1-\epsilon)\|y\|^{2}\leq\|\Pi y\|^{2}\leq(1+\epsilon)\|y\|^{2}
\]
can be rewritten as
\[
(1-\epsilon)U^{\top}U\preceq U^{\top}\Pi^{\top}\Pi U\preceq(1+\epsilon)U^{\top}U.
\]
Using $U^{\top}U=I_{d}$, we have
\[
(1-\epsilon)I_{d}\preceq U^{\top}\Pi^{\top}\Pi U\preceq(1+\epsilon)I_{d}.
\]
\end{proof}
For the special case $d=1$, the definition becomes
\[
\P\left[|\|\Pi a\|^{2}_{2}-1|\leq\epsilon\right]\geq1-\delta
\]
for all unit vectors $a$. The classical Johnson-Lindenstrauss lemma
is usually stated for finite point sets under random projection. The
OSE formulation asks instead for simultaneous preservation of every
vector in a fixed $d$-dimensional subspace; Gaussian, Bernoulli,
and more general random matrices can achieve this \cite{Vempala04}.
\begin{lem}[Johnson-Lindenstrauss Lemma]
Let $\Pi\in\R^{m\times n}$ be a random matrix with i.i.d. entries
from $N(0,\frac{1}{m})$ or uniformly sampled from $\pm\frac{1}{\sqrt{m}}$
with $m=\Theta(\frac{1}{\epsilon^{2}}\log(\frac{1}{\delta}))$. Then,
$\Pi$ is a $(1,\epsilon,\delta)$ OSE.
\end{lem}

We will skip the proof for this as we will prove a more general result
later. Next, we show that any OSE for $d=1$ is an OSE for general
$d$. Therefore, it suffices to focus on the case $d=1$. First, we
need a lemma about $\epsilon$-net on $S^{n-1}$.
\begin{lem}
For any $\epsilon>0$ and any $n\in\mathbb{N}$, there are at most
$(1+\frac{2}{\epsilon})^{n}$ unit vectors $x_{i}\in\Rn$ such that
for any unit vector $x\in\Rn$, there is an $i$ such that $\|x-x_{i}\|_{2}\leq\epsilon$.
\end{lem}

\begin{rem*}
We call the points $\{x_{i}\}_{i}$ the $\epsilon$-net on $S^{n-1}$.
\end{rem*}
\begin{proof}
We consider the following algorithm:
\begin{enumerate}
\item $V=S^{n-1}$. $i\leftarrow1$.
\item While there is $x_{i}\in V$
\begin{enumerate}
\item $V\leftarrow V\backslash B(x_{i},\epsilon)$.
\item $i\leftarrow i+1$.
\end{enumerate}
\end{enumerate}
Let $\{x_{1},x_{2},\cdots,x_{N}\}$ be the points it found. By the
construction, it is clear that $\|x-x_{i}\|\leq\epsilon$ for some
$i$ (else the procedure would continue). Now consider balls of radius
$\epsilon/2$ centered at each $x_{i}$. To bound the number of points,
we note that all $B(x_{i},\frac{\epsilon}{2})$ are disjoint and that
all balls lie in $B(0,1+\frac{\epsilon}{2})$. Therefore,
\[
N\cdot\vol(B(0,\frac{\epsilon}{2}))\leq\vol(B(0,1+\frac{\epsilon}{2})).
\]
and hence $N\leq(1+\frac{2}{\epsilon})^{n}$.
\end{proof}
\begin{lem}
\label{lem:OSE_reduction}Suppose that $\Pi$ is a $(1,\epsilon,\delta)$
OSE. Then, $\Pi$ is a $(d,4\epsilon,5^{d}\delta)$ OSE.
\end{lem}

\begin{proof}
Let $\mathcal{N}=\{x_{i}\}^{N}_{i=1}$ be a $\frac{1}{2}$-net for
$S^{d-1}$, with $N\leq5^{d}$. Put $\Delta=U^{\top}\Pi^{\top}\Pi U-I_{d}$.
By the $(1,\epsilon,\delta)$ guarantee for $\Pi$ and a union bound
over $U\mathcal{N}$, with probability at least $1-5^{d}\delta$,
\[
|y^{\top}\Delta y|=\left|\|\Pi Uy\|^{2}-1\right|\leq\epsilon\quad\text{for all }y\in\mathcal{N}.
\]
On this event, let $x$ be a unit eigenvector of $\Delta$ with $|x^{\top}\Delta x|=\|\Delta\|_{\op}$,
and choose $y\in\mathcal{N}$ with $\|x-y\|_{2}\leq\frac{1}{2}$.
Write $y=\alpha x+z$ with $z\perp x$. Then $\alpha\geq\frac{7}{8}$,
so
\[
|y^{\top}\Delta y|\geq(2\alpha^{2}-1)\|\Delta\|_{\op}\geq\frac{17}{32}\|\Delta\|_{\op}.
\]
Thus $\|\Delta\|_{\op}\leq\frac{32}{17}\epsilon\leq4\epsilon$, and
hence
\[
\P\left(\|U^{\top}\Pi^{\top}\Pi U-I_{d}\|_{\op}\leq4\epsilon\right)\geq1-5^{d}\delta.
\]
\end{proof}
This reduction and the Johnson-Lindenstrauss lemma show that a random
matrix with independent uniformly random entries in $\pm\frac{1}{\sqrt{m}}$
is a $(d,\epsilon,\delta)$ OSE with
\[
m=\Theta(\frac{1}{\epsilon^{2}}(d+\log(\frac{1}{\delta}))).
\]
As we discussed before any $(d,\epsilon,\delta)$ OSE should have
at least $d$ rows. Therefore, the number of rows of this OSE is tight
for the regime $\epsilon=\Theta(1)$. We only need an OSE for $\epsilon=\Theta(1)$
because of Lemma \ref{lem:Richardson}; by iterating we can get any
$\epsilon$ with an overhead of $\log(1/\epsilon)$. Unfortunately,
computing $\Pi A$ is in fact more expensive than $A^{\top}A$. The
first involves multiplying $\Theta(d)\times n$ and $n\times d$ matrix,
the second one involves multiplying $d\times n$ and $n\times d$
matrix.

\subsection{Sparse Johnson-Lindenstrauss}

We consider the sparse matrix $\Pi\in\R^{m\times n}$ where each entry
is $+\frac{1}{\sqrt{s}}$ with probability $\frac{s}{2m}$, $-\frac{1}{\sqrt{s}}$
with probability $\frac{s}{2m}$, and $0$ otherwise. We will show
that this random matrix is a $(d,\epsilon,\delta)$-OSE for $s=\Theta(\log^{2}(d/\delta)/\epsilon^{2})$
and $m=\Theta(d\log(\frac{d}{\delta})/\epsilon^{2})$. When $n=d=1$,
$\|\Pi x\|^{2}$ is $1/s$ times the number of nonzero entries in
its single column. Therefore, we indeed need that $s$ scales like
$1/\epsilon^{2}$. The advantage of a sparse embedding is that applying
it can be much more efficient.
\begin{rem}
It turns out that one can select exactly $s$ non-zeros for each column.
This allows us to use $s=\Theta(\frac{\log(d/\delta)}{\epsilon})$.
The proof of this is slightly more complicated due to the lack of
independence \cite{cohen2016nearly}.
\end{rem}

To analyze $U^{\top}\Pi^{\top}\Pi U$, we note that 
\[
U^{\top}\Pi^{\top}\Pi U=\sum^{m}_{r=1}(\Pi U)^{\top}_{r}(\Pi U)_{r}.
\]
Since each row of $\Pi$ is independent, we can use matrix concentration
bounds to analyze the sum above. See \cite{tropp2015introduction}
for a survey on matrix concentration bounds.
\begin{thm}[Matrix Chernoff]
\label{thm:matrix-chernoff}Suppose we have a sequence of independent,
random, self-adjoint matrices $M_{j}\in\R^{n\times n}$ such that
$\E M_{j}=I$ and $0\preceq M_{j}\preceq R\cdot I$. Then, for $T=\frac{R}{\varepsilon^{2}}\log\frac{n}{\delta}$,
\[
(1-O(\varepsilon))I\preceq\frac{1}{T}\sum^{T}_{j=1}M_{j}\preceq(1+O(\varepsilon))I
\]
with probability at least $1-\delta$.
\end{thm}

\begin{lem}[Hanson-Wright Inequality]
For independent random variables $\sigma_{1},\cdots,\sigma_{n}$
with $\E\sigma_{i}=0$, $|\sigma_{i}|\leq1$ and $A\in\R^{n\times n}$,
we have 
\[
\left|\sigma^{\top}A\sigma-\E\sigma^{\top}A\sigma\right|\lesssim\|A\|_{F}\sqrt{\log(\frac{1}{\delta})}+\|A\|_{\op}\log(\frac{1}{\delta})
\]
with probability $1-\delta$.
\end{lem}

For our problem, we set $M_{r}=m\cdot U^{\top}\pi_{r}\pi^{\top}_{r}U$
where $\pi_{r}$ is the $r$-th row of $\Pi$ (as a column vector).
One can check that $M_{r}\succeq0$ and that $\E M_{r}=I$. Next,
we note that
\begin{equation}
M_{r}\preceq m\cdot\pi^{\top}_{r}UU^{\top}\pi_{r}\cdot I.\label{matrix_ub}
\end{equation}
With small probability, $\pi^{\top}_{r}UU^{\top}\pi_{r}$ can be huge.
However, as long as $\pi^{\top}_{r}UU^{\top}\pi_{r}$ is bounded by
$R$ with probability $1-\delta$, then we can still use the matrix
Chernoff bound above. To bound $\pi^{\top}_{r}UU^{\top}\pi_{r}$,
we will use the following large deviation inequality.

\begin{lem}
Assume that $s\gg\frac{m}{n}\log(\frac{1}{\delta}),\quad s\gg\log^{2}(\frac{d}{\delta})/\epsilon^{2}$
and $m\gg d\log(\frac{d}{\delta})/\epsilon^{2}$. Then, 
\[
\pi^{\top}_{r}UU^{\top}\pi_{r}\leq\frac{\epsilon^{2}}{\log(d/\delta)}
\]
with probability $1-\delta$. (Here $\gg$ means greater by a sufficiently
large constant factor).
\end{lem}

\begin{proof}
Note that $\pi_{r}\in\Rn$ with each entry non-zero with probability
$\frac{s}{m}$. By the Chernoff bound, one can show that $\pi_{r}$
has at most $2sn/m$ non-zeros with probability $1-\delta$ using
that $sn/m\gg\log(\frac{1}{\delta})$. Let $I$ be the set of indices
for which $\pi_{r}$ is non-zero. Conditional on the set $I$ and
that $|I|\leq2sn/m$, note that $\sigma\defeq\pi_{r}|_{I}$ is a random
$\pm\frac{1}{\sqrt{s}}$ vector. Let $P=(UU^{\top})|_{I\times I}$.
Then, we have that
\[
\pi^{\top}_{r}UU^{\top}\pi_{r}=\sigma^{\top}P\sigma.
\]
The Hanson-Wright inequality shows that
\[
\left|\sigma^{\top}P\sigma-\E\sigma^{\top}P\sigma\right|\lesssim\frac{1}{s}\|P\|_{F}\sqrt{\log(\frac{1}{\delta})}+\frac{1}{s}\|P\|_{\op}\log(\frac{1}{\delta}).
\]
Note that $UU^{\top}$ is a projection matrix and hence $\|P\|_{\op}\leq1$.
Since $P\succeq0$, we have that
\[
\|P\|_{F}=\sqrt{\tr P^{2}}\leq\sqrt{\|P\|_{\op}\cdot\tr P}\leq\sqrt{\tr P}.
\]
Also, we have that $\E\sigma^{\top}P\sigma=\frac{1}{s}\tr P$. Hence,
we have
\[
\sigma^{\top}P\sigma\leq\frac{1}{s}\left(\tr P+O\left(\sqrt{\tr P\cdot\log\frac{1}{\delta}}+\log\frac{1}{\delta}\right)\right)
\]
with probability $1-\delta$. Note that $\tr UU^{\top}=\tr U^{\top}U=d$
and that $P$ is a principal submatrix of $UU^{\top}$ of size at
most $2sn/m$. By the Chernoff bound, one can show that
\[
\tr P\leq\frac{4sd}{m}.
\]
Hence, we have
\[
\sigma^{\top}P\sigma\leq\frac{4d}{m}+\frac{1}{s}O\left(\sqrt{\frac{sd}{m}\cdot\log\frac{1}{\delta}}+\log\frac{1}{\delta}\right)
\]
with probability $1-\delta$. Using $s\gg\log^{2}(\frac{d}{\delta})/\epsilon^{2}$
and $m\gg d\log(\frac{d}{\delta})/\epsilon^{2}$, we have $\sigma^{\top}P\sigma\leq\frac{\epsilon^{2}}{\log(d/\delta)}$
with probability $1-\delta$.
\end{proof}
Now we can prove the main theorem.
\begin{thm}
\label{thm:sparse_OSE}Consider a random sparse matrix $\Pi\in\R^{m\times n}$
where each entry is $+\frac{1}{\sqrt{s}}$ with probability $\frac{s}{2m}$,
$-\frac{1}{\sqrt{s}}$ with probability $\frac{s}{2m}$, and $0$
otherwise. There exist constants $c_{1},c_{2}$ such that for $m=c_{1}\frac{d\log(\frac{d}{\delta})}{\epsilon^{2}}$
and $s=c_{2}\frac{\log^{2}(\frac{d}{\delta})}{\epsilon^{2}}$ $\Pi$
is a $(d,\epsilon,\delta)$-OSE.
\end{thm}

\begin{proof}
Apply the previous lemma with failure probability $\delta/(2m)$.
By a union bound, with probability at least $1-\delta/2$, every row
satisfies
\[
\pi^{\top}_{r}UU^{\top}\pi_{r}\leq\frac{\epsilon^{2}}{\log(2d/\delta)}.
\]
On this event, using Theorem \ref{thm:matrix-chernoff} and the bound
(\ref{matrix_ub}), and choosing the constants in $m$ and $s$ large
enough, we have that
\[
(1-O(\varepsilon))I\preceq U^{\top}\Pi^{\top}\Pi U\preceq(1+O(\varepsilon))I.
\]
This is exactly the OSE guarantee for the subspace spanned by $U$.
\end{proof}

\subsection{Back to Regression}

In the last subsection, we presented a proof sketch of a suboptimal
OSE. For completeness, we state a tighter version here:

\begin{algorithm2e}[H]

\caption{$\mathtt{RegressionUsingOSE}$}

\SetAlgoLined

\textbf{Input:} a matrix $A\in\R^{n\times d}$, a vector $b\in\R^{n}$,
success probability $\delta$ and target accuracy $\epsilon$.

Let $\Pi$ be a $(d,\frac{1}{4},\delta)$ OSE constructed by Theorem
\ref{thm:sparse_OSE}.

Compute $A'=2\Pi A$, $M=A'^{\top}A'$ and $M^{-1}$.

$x^{(1)}=0$.

\For{$k=1,2,\cdots,4\log(\frac{1}{\epsilon})$}{

$x^{(k+1)}=x^{(k)}-M^{-1}(A^{\top}Ax^{(k)}-A^{\top}b)$

}

\Return $x^{\last}$.

\end{algorithm2e}
\begin{thm}
Given any full-column-rank matrix $A\in\R^{n\times d}$, any vector
$b\in\R^{n}$, the algorithm $\mathtt{RegressionUsingOSE}$ returns
a vector $x$ such that 
\[
\|A^{\top}Ax-A^{\top}b\|_{(A^{\top}A)^{-1}}\leq\epsilon\|A^{\top}b\|_{(A^{\top}A)^{-1}}
\]
with probability $1-\delta$ in $\widetilde{O}(\nnz(A)+d^{\omega})$
time.
\end{thm}

\begin{proof}
The guarantee of $x$ follows from the definition of OSE and Lemma
\ref{lem:Richardson}. We note that $\Pi A$ simply involves duplicating
each row of $A$ into roughly $s$ many rows in $\Pi A$. Hence, the
cost of computing $\Pi A$ is $O(s\cdot\nnz(A))=\widetilde{O}(\nnz(A))$.
The cost of computing $M$ takes $\widetilde{O}(d^{\omega})$ time.
Computing $M^{-1}$ takes $\widetilde{O}(d^{\omega})$ time. The loop
takes $\widetilde{O}(\nnz(A))$ time. This explains the total time.
\end{proof}
Linear regression can also be solved in time $O\left(\nnz(A)+d^{O(1)}\right)$
via a very sparse embedding: each column of $\Pi$ picks exactly one
nonzero entry in a random location. This was analyzed by Clarkson
and Woodruff \cite{woodruff2014sketching}. See also \cite{kannan_vempala_2017}
for a survey.
